\exercisesforappendix

\begin{enumerate}
\def\labelenumi{\arabic{enumi}.}
\tightlist
\item
  令 \(\Phi_n\) 为指标为 \(n\) 的圆分多项式（《代数》，第五章，§ 11，第 2 号）。利用公式
\end{enumerate}

\[
\mathbf {X} ^ {n} - 1 = \prod_ {d | n} \Phi_ {d} (\mathbf {X}),
\]

证明

\[
\Phi_ {n} (\mathrm{X}) = \prod_ {d | n} \left(\mathrm{X} ^ {n / d} - 1\right) ^ {\mu (d)}.
\]

（将 Möbius 反演公式应用于域 \(\mathbf{Q}(\mathbf{X})\) 的乘法群。）

\begin{enumerate}
\def\labelenumi{\arabic{enumi}.}
\setcounter{enumi}{1}
\tightlist
\item
  令 D 为幺半群 \(\mathbf{N}^*\) 的全代数（《代数》，第三章，§ 2，第 10 号）。若 \(n \in \mathbf{N}^*\)，令 \(n^\omega\) 表示它在 D 中的像，于是 \(1^\omega = 1\)，且 \((nm)^\omega = n^\omega m^\omega\) 对 \(n\)、\(m\) 属于 \(\mathbf{N}^*\) 成立。D 中每个元素 \(f\) 都能唯一地写成级数 \(\sum_{i=1}^{\infty}\)
\end{enumerate}

\(\sum_{n=1}^{\infty} a_n n^\omega\)，其中 \(a_n \in \mathbf{K}\)。

\begin{enumerate}
\def\labelenumi{(\alph{enumi})}
\item
  令 \(f = \sum_{n=1}^{\infty} a_n n^\omega\) 是 D 中的元素。证明 \(f\) 在 D 中可逆，当且仅当 \(a_1\) 在 K 中可逆。特别地，若 K 是局部环，则 D 也是局部环。
\item
  记 \(\zeta = \sum_{n=1}^{\infty} n^{\omega}\)，\(\mu = \sum_{n=1}^{\infty} \mu(n)n^{\omega}\)。证明 \(\zeta\) 与 \(\mu\) 互为逆元。推出若 \(s = \sum_{n} s(n)n^{\omega}\) 和 \(t = \sum_{n} t(n)n^{\omega}\) 是 D 中的两个元素，则关系 \(s = \zeta\) . \(t\) 与 \(t = \mu\) . \(s\) 等价（Möbius 反演公式的一个变体）。
\item
  令 \(\mathbf{P}\) 为素数集合。证明族 \((1 - p^{\omega})\) 对 \(p\in \mathbf{P}\) 可在 \(\mathbf{D}\) 中相乘，并且
\end{enumerate}

\[
\mu = \sum_ {p \in P} ^ {\infty} (1 - p ^ {\omega}) \quad \text { and } \quad \zeta = \prod_ {p \in P} \frac {1}{1 - p ^ {\omega}}.
\]

